\chapter{Конструкторская часть}

\section{Требования к программному обеспечению}

К разрабатываемой программе предъявлен ряд требований:

\textbf{Входные данные:} Две матрицы

\textbf{Выходные данные:} Матрица, являющаяся результатом умножения данных матриц

\begin{itemize}
    \item Должна быть возможность одиночного выполнения умножения введенных пользователем матриц введенных пользователем 4 размеров;
    \item Должна быть возможность массированного замера времени выполнения реализаций алгоритмов умножения матриц;
    \item Должна быть возможность замера процессорного времени программы;
    \item Должна быть возможность вывода графиков и таблиц замеров процессорного времени.
\end{itemize}

\section{Разработка алгоритмов}

Для всех алгоритмов заданы:
\begin{itemize}
    \item $rows\_A$ -- количество строк в $A$;
    \item $cols\_A$ -- количество столбцов в $A$;
    \item $rows\_B$ -- количество строк в $B$;
    \item $cols\_B$ -- количество столбцов в $B$;
    \item Матрица $result$ размером $rows\_A$ на $cols\_B$, заполненная нулями.
\end{itemize}

Для алгоритмов Винограда заданы:
\begin{itemize}
    \item Массив $mulH$ размера M, заполненный нулями;
    \item Массив $mulV$ размера Q, заполненный нулями;
\end{itemize}

Cтандартный алгоритм умножения матриц представлен в листинге \ref{alg:std_des}.

\begin{algorithm}[H]
\caption{Cтандартный алгоритм умножения матриц}
\label{alg:std_des}
\begin{algorithmic}
    \For{$i = 0$ \textbf{to} $rows\_A - 1$}
        \For{$j = 0$ \textbf{to} $cols\_B - 1$}
            \For{$k = 0$ \textbf{to} $cols\_A - 1$}
                \State $result[i][j] = result[i][j] + A[i][k] \times B[k][j]$
            \EndFor
        \EndFor
    \EndFor
    
    \State \textbf{Return} $result$
\end{algorithmic}
\end{algorithm}

Алгоритм умножения матриц Винограда представлен в листинге \ref{alg:vin_des}.

\begin{algorithm}[H]
\caption{Алгоритм умножения матриц Винограда}
\label{alg:vin_des}
\begin{algorithmic}
    \For{$i = 0$ \textbf{to} $rows\_A - 1$}
        \For{$j = 0$ \textbf{to} $cols\_A // 2 - 1$}
            \State $mulH[i] = mulH[i] + A[i][2 \cdot j] \times A[i][2 \cdot j + 1]$
        \EndFor
    \EndFor
    \For{$i = 0$ \textbf{to} $cols\_B - 1$}
        \For{$j = 0$ \textbf{to} $rows\_B // 2 - 1$}
            \State $mulV[i] = mulV[i] + B[2 \cdot j][i] \times B[2 \cdot j + 1][i]$
        \EndFor
    \EndFor
    
    \For{$i = 0$ \textbf{to} $rows\_A - 1$}
        \For{$j = 0$ \textbf{to} $cols\_B - 1$}
            \State $result[i][j] = -mulH[i] - mulV[j]$
            \For{$k = 0$ \textbf{to} $cols\_A // 2 - 1$}
                \State $result[i][j] = result[i][j] + (A[i][2 \cdot k] + B[2 \cdot k + 1][j]) \times (A[i][2 \cdot k + 1] + B[2 \cdot k][j])$
            \EndFor
        \EndFor
    \EndFor
    
    \If{$cols\_A \mod 2 == 1$}
        \For{$i = 0$ \textbf{to} $rows\_A - 1$}
            \For{$j = 0$ \textbf{to} $cols\_B - 1$}
                \State $result[i][j] = result[i][j] + A[i][cols\_A - 1] \times B[cols\_A - 1][j]$
            \EndFor
        \EndFor
    \EndIf
    \State \textbf{Return} $result$
\end{algorithmic}
\end{algorithm}

Оптимизированный алгоритм умножения матриц Винограда представлен в листинге \ref{alg:opt_vin_des}.(Двоичный сдвиг вместо умножения на 2, введение декремента при вычислении вспомогательных массивов, выделение начальной итерации из каждого внешнего цикла)

\begin{algorithm}[H]
\caption{Оптимизированный алгоритм умножения матриц Винограда}
\label{alg:opt_vin_des}
\begin{algorithmic}
    \For{$i = 0$ \textbf{to} $rows\_A - 1$}
        \State $mulH[i] = -A[i][0] \times A[i][1]$
        \For{$j = 1$ \textbf{to} $cols\_A // 2 - 1$}
            \State $mulH[i] = mulH[i] - A[i][j << 1] \times A[i][(j << 1) + 1]$
        \EndFor
    \EndFor
    \For{$i = 0$ \textbf{to} $cols\_B - 1$}
        \State $mulV[i] = -B[0][i] \times B[1][i]$
        \For{$j = 1$ \textbf{to} $rows\_B // 2 - 1$}
            \State $mulV[i] = mulV[i] - B[j << 1][i] \times B[(j << 1) + 1][i]$
        \EndFor
    \EndFor
    
    \For{$i = 0$ \textbf{to} $rows\_A - 1$}
        \For{$j = 0$ \textbf{to} $cols\_B - 1$}
            \State $result[i][j] = mulH[i] + mulV[j]$
            \State $result[i][j] = result[i][j] + (A[i][0] + B[1][j]) \times (A[i][1] + B[0][j])$
            \For{$k = 1$ \textbf{to} $cols\_A // 2 - 1$}
                \State $result[i][j] = result[i][j] + (A[i][k << 1] + B[(k << 1) + 1][j]) \times (A[i][(k << 1) + 1] + B[k << 1][j])$
            \EndFor
        \EndFor
    \EndFor
    
    \If{$cols\_A \mod 2 == 1$}
        \For{$i = 0$ \textbf{to} $rows\_A - 1$}
            \For{$j = 0$ \textbf{to} $cols\_B - 1$}
                \State $result[i][j] = result[i][j] + A[i][cols\_A - 1] \times B[cols\_A - 1][j]$
            \EndFor
        \EndFor
    \EndIf
    \State \textbf{Return} $result$
\end{algorithmic}
\end{algorithm}

\section{Модель вычислений}

В данной работе будет выполняться оценка трудоемкости алгоритмов, поэтому будет введена модель вычислений:

\textbf{Трудоемкость базовых операций}

\begin{itemize}
    \item Единичные операции: $=, ==, -, +, <<, []$
    \item Операции трудоемкости 2: $\times, \mod, \cdot, //$
\end{itemize}

\textbf{Трудоемкость операторов цикла}

Пусть для цикла вида:
for инициализация to значение сравнения do тело
    
Известны трудоемкости соответствующих блоков $f_{\text{инициализация}}$,\newline$f_{\text{значение сравнения}}$, $f_{\text{тело}}$, тогда трудоемкость цикла имеет следующий вид:

\begin{equation}
    \label{eq:work1}
    f = f_{\text{инициализация}} + f_{\text{значение сравнения}} + 1 + M*(f_{\text{тело}} + 1 + f_{\text{значение сравнения}} + 1)
\end{equation}

\textbf{Трудоемкость условного оператора}

Пусть трудоемкость условного перехода равна 0. При известных $f_\text{усл}$, $f_{\text{тело}}$, тогда трудоемкость цикла вида:

if усл then тело end if

Трудоемкость равна:

\begin{equation}
    \label{eq:work2}
    f = f_{\text{усл}} + [0\text{, если лучший случай, иначе } f_{\text{тело}}]
\end{equation}

\section{Оценка трудоемкости}

Пусть обрабатываются матрицы размера $M \times N$ и $N \times Q$

\textbf{Cтандартный алгоритм умножения матриц}

Лучший и худший случай идентичны, так как алгоритм всегда выполняет одно и то же количество операций, независимо от структуры матриц.

\begin{equation}
    \begin{aligned}
    f &= \underbrace{3 + M \times (\underbrace{3 + N \times (\underbrace{3 + Q \times (12 + 3)}_{\text{цикл по k}} + 3)}_{\text{цикл по j}} + 3)}_{\text{цикл по i}} = \\
    &= 3+6M+6MN+15MNQ
    \end{aligned}
\end{equation}

\textbf{Алгоритм умножения матриц Винограда}

Лучший случай - N mod 2 != 1:

\begin{equation}
    \begin{aligned}
    f &=
    \underbrace{3 + M \times (\underbrace{
    5 + N/2 \times (15+ 5)
    }_{цикл по j} + 3)}_{цикл по i} + \\
    &+ \underbrace{3 + Q \times (\underbrace{
    5 + N/2 \times (15+ 5)
    }_{цикл по j} + 3)}_{цикл по i} + \\
    &+ \underbrace{3 +M \times (\underbrace{3 + Q \times (7 + \underbrace{
    5 + N/2 \times (28 + 5)
    }_{цикл по k} + 3)}_{\text{цикл по j}} + 3)}_{\text{цикл по i}} +
    \underbrace{3}_{if} = \\
    &= 12+14M+8Q+10MN+15MQ+10NQ+16.5MNQ
    \end{aligned}
\end{equation}

Худший случай - N mod 2 == 1:

\begin{equation}
    \begin{aligned}
    f &=
    \underbrace{3 + M \times (\underbrace{
    5 + N/2 \times (15+ 5)
    }_{цикл по j} + 3)}_{цикл по i} + \\
    &+ \underbrace{3 + Q \times (\underbrace{
    5 + N/2 \times (15+ 5)
    }_{цикл по j} + 3)}_{цикл по i} + \\
    &+ \underbrace{3 + M \times (\underbrace{3 + Q \times (7 + \underbrace{
    5 + N/2 \times (28 + 5)
    }_{цикл по k} + 3)}_{\text{цикл по j}} + 3)}_{\text{цикл по i}} + \\
    &+ \underbrace{3 +
    \underbrace{3 + M \times (\underbrace{3 + Q \times (14 + 3)}_{\text{цикл по j}} + 3)}_{\text{цикл по i}}
    }_{if} = \\
    &= 15+20M+8Q+10MN+32MQ+10NQ+16.5MNQ
    \end{aligned}
\end{equation}

\textbf{Оптимизированный алгоритм умножения матриц Винограда}

Лучший случай - N mod 2 != 1:

\begin{equation}
    \begin{aligned}
    f &=
    \underbrace{3 + M \times (9 + \underbrace{
    5 + (N/2-1) \times (13 + 5)
    }_{цикл по j} + 3)}_{цикл по i} + \\
    &+ \underbrace{3 + Q \times (\underbrace{
    9 + 5 + (N/2-1) \times (13 + 5)
    }_{цикл по j} + 3)}_{цикл по i} + \\
    &+ \underbrace{3 + M \times (\underbrace{3 + Q \times (6 + 18 + \underbrace{
    5 + (N/2-1) \times (23 + 5)
    }_{цикл по k} + 3)}_{\text{цикл по j}} + 3)}_{\text{цикл по i}} +
    \underbrace{3}_{if} = \\
    &= 12+5M+9MN+4MQ+9MN+14MNQ
    \end{aligned}
\end{equation}

Худший случай - N mod 2 == 1:

\begin{equation}
    \begin{aligned}
    f &=
    \underbrace{3 + M \times (9 + \underbrace{
    5 + (N/2-1) \times (13 + 5)
    }_{цикл по j} + 3)}_{цикл по i} + \\
    &+ \underbrace{3 + Q \times (\underbrace{
    5 + (N/2-1) \times (13 + 5)
    }_{цикл по j} + 3)}_{цикл по i} + \\
    &+ \underbrace{3 + M \times (\underbrace{3 + Q \times (6 + 18 + \underbrace{
    5 + (N/2-1) \times (23 + 5)
    }_{цикл по k} + 3)}_{\text{цикл по j}} + 3)}_{\text{цикл по i}} + \\
    &+ \underbrace{3 + 
    \underbrace{3 + M \times (\underbrace{3 + Q \times (14 + 3)}_{\text{цикл по j}} + 3)}_{\text{цикл по i}}
    }_{if} = \\
    &= 15+11M+9MN+21MQ+9MN+14MNQ
    \end{aligned}
\end{equation}

\section{Вывод}

В результате конструкторской части были определены требования к ПО, спроектированы стандартный алгоритм умножения матриц, алгоритм Винограда и оптимизированный по варианту алгоритм Винограда, а также произведена оценка трудоемкости алгоритмов.
